\begin{description}
\item[a)] (B) the one on December 16 \hfill(1\,p)

  At a given moment of the day, the Sun elevation is lower in winter on the
  northern hemisphere.

\item[b)] (B) late in the morning \hfill(1\,p)

  The shadows point slightly west of north. In fact these satellite images are
  taken at around 9:45 UTC (10:45 Central European Time).

\item[c)] The measured shadow lengths are $x_1 = 125$\,m (epoch 1, summer
  solstice) and $x_2 = 780$\,m (epoch 2, winter solstice).

  Let the unknown height of the chimney be $h$. The chimney and its shadow form
  the legs of a right-angled triangle and thus
  \begin{align*}
    \tan\alpha_1 &= h/x_1 \tag{6.1}\\
    \tan\alpha_2 &= h/x_2, \tag{6.2}
  \end{align*}
  \hfill(1\,p)

  where $\alpha_1$ and $\alpha_2$ are the Sun elevation angles at epoch 1 and
  epoch 2, respectively. \hfill(1\,p)

  At the time of the summer and winter solstices, the Earth is located at
  opposite sides of the Sun during its orbital revolution. Since the Equator is
  tilted by $\varepsilon \approx 23.5^\circ$ to the orbital plane, at a given
  point of the surface, and at a given time of the day, the Sun elevation
  angles will differ by $\alpha_1 - \alpha_2 = 2\varepsilon$ (here we ignore
  refraction and that the satellite images were not taken on the exact solstice
  days). \hfill(4\,p)

  Notice that, to a good approximation, $\alpha_1 - \alpha_2 \approx 45^\circ$.
  Assuming $\alpha_1 = \alpha_2 + 45^\circ$ makes the solution a lot easier.
  Using the trigonometric identity
  \[
    \tan\alpha_1 = \tan(\alpha_2 + 45^\circ)
      = \frac{\tan\alpha_2 + \tan 45^\circ}{1 - \tan\alpha_2 \tan 45^\circ},
    \tag{6.3}
  \]
  and substituting $\tan 45^\circ = 1$, Eq.~(6.3) becomes
  \[
    \tan\alpha_1 = \frac{1 + \tan\alpha_2}{1 - \tan\alpha_2}. \tag{6.4}
  \]
  \hfill(4\,p)

  Substituting $\tan\alpha_1$ and $\tan\alpha_2$ from Eqs.~(6.1) and (6.2) into
  Eq.~(6.4), one gets a quadratic equation for the unknown $h$:
  \[
    \frac{h}{x_1} = \frac{1 + \dfrac{h}{x_2}}{1 - \dfrac{h}{x_2}}. \tag{6.5}
  \]
  \hfill(2\,p)

  Reducing this equation to the standard format
  \[
    h^2 + h(x_1 - x_2) + x_1 x_2 = 0, \tag{6.6}
  \]
  and according to the quadratic formula, the two roots are obtained as
  \[
    h_{1,2} = \frac{x_2 - x_1 \pm \sqrt{(x_1 - x_2)^2 - 4x_1x_2}}{2}. \tag{6.7}
  \]
  \hfill(1\,p)

  Now with the measured values of $x_1 = 125$\,m and $x_2 = 780$\,m, the two
  roots for $h$ are $h_1 = 426$\,m and $h_2 = 229$\,m. \hfill(1\,p)

  Obviously, only one of the solutions is valid for our case.

  One may simply notice that $h_1 \approx 426$\,m seems just too high for a
  chimney (would in fact be comparable to the Empire State Building in New
  York, with 443\,m; the Eiffel tower in Paris is only 324\,m high).

  Alternatively, substituting $h_1 = 426$\,m into Eq.~(6.2), we obtain the Sun
  elevation angle at around the winter solstice,
  $\alpha_2 = \arctan(h_1/x_2) \approx 28.7^\circ$. However, at the geographic
  latitudes of Hungary, the Sun does not rise above ${\approx}19^\circ$ at this
  time of the year, so the larger value obtained for $h$ is clearly invalid.
  \hfill(2\,p)

  Note that according to the official data, the Tiszaújváros power plant has a
  chimney as high as $h = 250$\,m.

\item[d)] (B) and (E) \hfill(2\,p)
\end{description}
